\begin{resumo}

A classificação de textos curtos em português enfrenta três obstáculos ao mesmo tempo. O vocabulário é abreviado e ruidoso, como nas descrições de produtos de notas fiscais eletrônicas. As categorias são centenas, e desbalanceadas. E o método depende de grandes volumes de rótulos, cuja obtenção é o principal custo desses sistemas. Esta tese propõe e avalia o FALCO (\textit{Framework} de Aprendizado Ativo com LLM para texto CurtO), que reduz esse custo em três etapas. A primeira é a composição informada do conjunto inicial de treinamento. A segunda é a seleção ativa de instâncias por incerteza. A terceira é a substituição parcial do anotador humano por oráculos baseados em modelos de linguagem de grande porte (LLMs), organizados em fases de custo crescente. A investigação usa um conjunto público auditado de 250.365 descrições de produtos em 621 categorias, das quais 250.221 na versão corrigida empregada nos experimentos. O protocolo é pré-registrado, com partições imutáveis, análise pareada (McNemar, Wilcoxon) e artefato rastreável para cada número reportado. Cinco resultados sustentam a tese. (i) A composição do conjunto inicial importa mais quanto menor o orçamento: só o sorteio produz amplitude de $6{,}4$ pontos percentuais de acurácia com $|L_0|=100$. A heurística determinística proposta, o DRI-SL, combina densidade semântica e diversidade lexical sem usar nenhum rótulo. Ela supera o melhor indivíduo de um algoritmo genético supervisionado em todos os tamanhos de 100 a 5.000. A comparação é conservadora, pois a correção de circularidade introduzida na tese reduziu o envelope evolutivo em até $6{,}3$ pontos percentuais. (ii) Oráculos LLM \textit{zero-shot} atingem $77$--$83\%$ de acurácia no espaço fechado de 621 categorias, a custos de US\$~0,035 a US\$~0,92 por mil rótulos. O Macro F1 deles supera o de um classificador supervisionado leve treinado com a base completa de 250 mil descrições rotuladas. Quase metade dos seus erros concentra-se em convenções de catálogo: categorias-irmãs e classes guarda-chuva. (iii) O instrumento de medição é parte do resultado. Sem restrição de saída ao espaço de classes, $6{,}8\%$ das respostas são contadas como falsos erros. Regras de fronteira no \textit{prompt} melhoram o modelo fraco na distribuição de produção ($+4{,}6$ p.p., $p=0{,}012$), ao custo de degradação severa nas classes raras ($p<0{,}001$). O provedor de \textit{serving} também é instrumento: o mesmo modelo gratuito diverge significativamente de si mesmo entre dois serviços ($p<0{,}001$). Servido adequadamente, porém, ele empata com oráculos pagos na amostra de produção. (iv) As estratégias de incerteza superam a seleção aleatória na significância máxima que 8 sementes permitem ($p=0{,}0078$), e recuperam $78\%$ do Macro F1 do teto supervisionado com $15\%$ do \textit{pool} daquele experimento. A vantagem sobrevive intacta a ruído uniforme de oráculo até $\varepsilon=0{,}4$, faixa que cobre o erro observado nos LLMs reais. Nenhum oráculo alcançou o critério pré-registrado de $85\%$ de acurácia. O gate metodológico definiu então a configuração final: oráculo econômico na fase inicial e oráculo de maior acurácia na fase avançada. Ele também tornou o aprendizado com rótulos ruidosos o cenário central. (v) O framework foi executado ponta a ponta com oráculo LLM gratuito, a custo monetário zero. A validação usou o classificador forte, o BERTimbau, com três sementes e avaliação na população reservada de 177.490 instâncias. Ela respondeu à hipótese central: obter pelo menos $95\%$ da acurácia da supervisão completa do \textit{pool} de referência, a métrica do pré-registro, com no máximo 34.724 rótulos ($15\%$ da população deduplicada). O veredito é duplo. \textbf{O critério é atingível dentro do teto, e é cruzado com 20 mil rótulos ($8{,}6\%$ da população) nas três sementes, em varredura \textit{post hoc} com rótulos de gabarito. A configuração executada, porém, parou por estagnação aquém do piso, com 11.936 rótulos ($5{,}2\%$)}. A decomposição pareada não atribui a perda ao ruído do oráculo, cujo custo se concentra na acurácia e não em Macro F1. Também não a atribui à seleção: a perda vem do critério de parada herdado do classificador leve. Em Macro F1, análise de robustez introduzida nesta tese, o critério é atingido com 30 mil rótulos ($13{,}0\%$ da população), também nas três sementes e dentro do teto. E, na média das três sementes, 35 mil rótulos ativamente selecionados superam a supervisão completa do \textit{pool} nas duas métricas, com heterogeneidade declarada: duas sementes mantêm vantagem significativa, uma inverte em Macro F1 e empata em acurácia. Os experimentos em escala populacional mostram ainda que a autoavaliação em dados coletados ativamente é enviesada, em direção que depende da métrica: a acurácia da autoavaliação subestima o desempenho real, e o Macro F1 da autoavaliação o superestima. Isso fundamenta o conjunto reservado de avaliação e o critério de liberação que o framework adota. As limitações de primeira ordem são declaradas. Os resultados provêm de uma única base, do próprio domínio de estudo. A superação da supervisão completa pelo braço de 35 mil rótulos vale na média das sementes, com uma semente divergente em Macro F1. E a análise de custo apoia-se em razões entre oráculos, como a diferença de 26 vezes sob empate estatístico de acurácia, não em preços absolutos, que datam. A tese entrega ainda a métrica LCE (\textit{Learning Curve Efficiency}), para comparação de curvas de aprendizado, a biblioteca aberta \texttt{activelearning} e a ferramenta FlowBuilder. Com ingestão saneada de dados, execução parametrizada e curvas de aprendizado, elas tornam o protocolo reutilizável em outros domínios.

\end{resumo}
